% =====================================================================
% SUPERSEDED 2026-08-29 -- DO NOT \input THIS FILE.
% Its content has been integrated into chapter3.tex as
% \subsubsection{Field reconstruction} (label subsubsec:field_recon).
% This copy is kept only as a record of the earlier draft; it predates
% the CORAL decoder result and the flap-separation finding, and it
% redefines tab:meansplit / fig:meansplit_decomp / fig:meansplit_fields,
% so including it now would raise duplicate-label errors.
% ---------------------------------------------------------------------
% Field reconstruction and the near-wall bias (mean-split study)
% ---------------------------------------------------------------------
% Drop-in subsubsection for Sec. 4.1 (subsec:training_results), replacing
% the commented placeholder "Network depth and field reconstruction".
% Requires in the preamble:  \newcommand{\msCcBaseNear}{0.140}
\newcommand{\msCcMsNear}{0.0244}
\newcommand{\msCcFactor}{5.7}
\newcommand{\msABaseNear}{0.194}
\newcommand{\msAMsNear}{0.0082}
\newcommand{\msAFactor}{24}

% Figures (Thesis_article_format/Chapters/Chapter 3/Images/):
%   fig_meansplit_decomp.png   fig_meansplit_fields.png
% Numbers are verified against analysis_meansplit/meansplit_arm_summary.csv
% (30 runs, 5 seeds/arm; see recon/analysis/MEANSPLIT_NOTES.md).
% =====================================================================

\subsubsection{Field reconstruction and the near-wall bias}
\label{sssec:field_recon}

The same latent-dynamics/reconstruction pair used for the loads can be queried over
the full spatial domain: replacing the sectional-force head with a three-output head
$(v_x,v_y,p)$ turns $\mathcal{NN}_\mathrm{rec}(z(t),u(t),x,y)$ into a coordinate
decoder of the flow field, trained by teacher forcing against the URANS snapshots.
This exercises the surrogate far more severely than the integrated loads, and exposes
where its representation is weakest. Measuring the reconstruction error per field and
per spatial region\,---\,near-airfoil, near-wake, far-field and the airfoil surface,
each normalised by that field's range\,---\,the $x$-velocity error concentrates in the
boundary layer and near wake, reaching \msCcBaseNear\ (gust\,+\,flap) and \msABaseNear\
(gust only) NRMSE in the near-airfoil band, several times the far-field value that the
global metric is dominated by. Decomposing this error into a time-mean (static) part
and a fluctuating (dynamic) part reveals that it is almost entirely a persistent
mean-flow bias: the static component accounts for $97\%$ of the near-airfoil $v_x$
error (red bars, \cref{fig:meansplit_decomp}). The plain tanh coordinate decoder,
with raw $(x,y)$ inputs, smears the sharp mean boundary-layer gradient and never
recovers it, while tracking the gust/flap fluctuations comparatively well.

\paragraph{Mean-field split.} Because the offending component is static, it need not be
learned at all. We store the training-set mean field $\bar{u}(x,y)$ once and train the
decoder on the fluctuations $u-\bar{u}$, recomputing the min--max normalisation on the
fluctuation ranges and adding $\bar{u}$ back at inference; the reported errors remain in
total-field units and directly comparable to the baseline. This is the cheapest possible
change\,---\,no new inputs, architecture or loss terms, and the latent dynamics are
untouched\,---\,and it mirrors the mean/POD-mode split standard in projection-based
reduced-order models~\cite{BennerGugercinWillcox2015}. The near-airfoil $v_x$ error drops to
\msCcMsNear\ (a factor \msCcFactor) on the gust\,+\,flap case and \msAMsNear\ (a factor
\msAFactor) on the gust-only case; the static share of the near-wall error collapses
from $97\%$ to $3\%$ (blue bars, \cref{fig:meansplit_decomp}), while the dynamic
component is left unchanged\,---\,the split removes the bias without disturbing what the
network already captured. \Cref{tab:meansplit} lists the region-resolved figures, and
\cref{fig:meansplit_fields} shows the effect in physical space: the error halo that
surrounds the entire profile in the baseline is all but eliminated, leaving only a
localised residual at the flap gap. The global combined metric moves less (from
$1.4\cdot10^{-2}$ to $7.2\cdot10^{-3}$) because it is dominated by the already-easy
far field and mixes the three physical units, a known artefact of the range-normalised
combined NRMSE; the per-field, per-region view is the informative one.

\paragraph{What the residual is.} The error that survives is dynamic, not static, and it
sits on the airfoil surface where the flap deflects (\cref{fig:meansplit_fields}, right).
A short ablation confirms that this floor is not reachable through the decoder inputs:
referencing the mean to the trim state instead of the ensemble mean, adding a
Tikhonov penalty, and appending wall-distance and boundary-layer-mask input
features\,---\,the standard remedy for near-wall gradients in steady
surrogates~\cite{Catalani2025}\,---\,each match but none beat the plain split beyond
seed scatter, and the wall features help only on the gust-only case, where the geometry
is static, giving no gain once the flap moves. Regularisation does not alter the
$d_s=1$ selection of \cref{subsec:training_results} for the field either. The remaining
near-wall error is therefore a property of the fluctuation dynamics captured by the
low-dimensional latent state, not of the spatial decoder, and would be addressed on the
dynamics side rather than by further input engineering.

% --------------------------------------------------------------------- table
\begin{table}[htbp]
\centering
\caption{Field-reconstruction $v_x$ error (NRMSE, normalised by the field range) by
spatial region, baseline coordinate decoder versus the mean-field split, on the two
test trajectories ($d_s=1$; mean over five seeds). ``static'' is the fraction of the
error carried by the time-mean bias.}
\label{tab:meansplit}
\begin{tabular}{llccc}
\toprule
Trajectory & Region & Baseline & Mean-split & Factor \\
\midrule
gust only
 & near-airfoil    & $0.194$ & $0.0082$ & $24\times$ \\
 & near-wake       & $0.0091$ & $0.0028$ & $3.2\times$ \\
 & far-field       & $0.0852$ & $0.0027$ & $31\times$ \\
 & airfoil surface & $0.0420$ & $0.0112$ & $3.7\times$ \\
\midrule
gust\,+\,flap
 & near-airfoil    & $0.140$ & $0.0244$ & $5.7\times$ \\
 & near-wake       & $0.0104$ & $0.0074$ & $1.4\times$ \\
 & far-field       & $0.0614$ & $0.0058$ & $11\times$ \\
 & airfoil surface & $0.0384$ & $0.0282$ & $1.4\times$ \\
\bottomrule
\end{tabular}
\end{table}

% --------------------------------------------------------------------- figures
\begin{figure}[htbp]
\centering
\includegraphics[width=\textwidth]{Thesis_article_format/Chapters/Chapter 3/Images/fig_meansplit_decomp.png}
\caption{Static (solid) versus dynamic (hatched) $x$-velocity reconstruction error by
region, for the baseline coordinate decoder (red) and the mean-field split (blue), on
the gust-only (left) and gust\,+\,flap (right) test trajectories. The baseline error is
an almost purely static near-wall bias; the split removes it while leaving the dynamic
component unchanged.}
\label{fig:meansplit_decomp}
\end{figure}

\begin{figure}[htbp]
\centering
\includegraphics[width=\textwidth]{Thesis_article_format/Chapters/Chapter 3/Images/fig_meansplit_fields.png}
\caption{Near-airfoil $x$-velocity at the gust peak of the gust\,+\,flap trajectory:
reference field (left) and absolute reconstruction error for the baseline (centre) and
the mean-field split (right), on a common colour scale. The baseline error surrounds the
whole profile; after the split only a localised residual at the flap gap remains.}
\label{fig:meansplit_fields}
\end{figure}
